\documentclass[11pt]{article}
\usepackage[a4paper,margin=25mm]{geometry}
\usepackage[T1]{fontenc}
\usepackage{lmodern}
\usepackage{amsmath,amssymb,amsthm,mathtools,microtype}
\usepackage{enumitem,xurl}
\usepackage[colorlinks=true,linkcolor=blue,citecolor=blue,urlcolor=blue]{hyperref}
\usepackage{authblk}
\renewcommand{\Affilfont}{\small}
\setlength{\parindent}{0pt}
\setlength{\parskip}{5pt}
\setlength{\emergencystretch}{2em}
\numberwithin{equation}{section}
\newtheorem{theorem}{Theorem}
\newtheorem{lemma}[theorem]{Lemma}
\newtheorem{proposition}[theorem]{Proposition}
\newtheorem{corollary}[theorem]{Corollary}
\theoremstyle{remark}
\newtheorem{remark}[theorem]{Remark}
\newcommand{\Tr}{\operatorname{Tr}}
\newcommand{\one}{\mathbf 1}
\newcommand{\R}{\mathbb R}
\newcommand{\norm}[1]{\lVert #1\rVert}
\newcommand{\ip}[2]{\langle #1,#2\rangle}
\newcommand{\pos}[1]{(#1)_+}
\title{Goodman-type lower bounds for odd cycles II: The intermediate regime}
\author[1,2]{Taeyoung~Kim}
\author[3]{Jineon~Baek}
\author[4]{Seonghyuk~Im}
\author[2]{Joonkyung~Lee}
\author[1,2]{Hongseok~Yang}
\affil[1]{School of Computing, KAIST, Korea}
\affil[2]{School of Computational Sciences, Korea Institute for Advanced Study (KIAS), Korea}
\affil[3]{June E Huh Center for Mathematical Challenges, Korea Institute for Advanced Study (KIAS), Korea}
\affil[4]{Center for AI and Natural Science, Korea Institute for Advanced Study (KIAS), Korea}
\affil[ ]{\newline\texttt{taeyoungkim21@kaist.ac.kr}, \texttt{{jineon.kias}@gmail.com}, \texttt{seonghyuk@kias.re.kr}, \texttt{joonkyunglee@kias.re.kr}, \texttt{hongseokyang@kias.re.kr}}
\date{September 23, 2026}
\begin{document}
\maketitle
\begin{abstract}
We give a self-contained proof of
\[
 t(C_m,W)\ge p^m-p(1-p)^{m-1}
\]
for every graphon of edge density $0\le p\le2/3$ and every odd $m\ge3$.
The nontrivial range is $1/2<p\le2/3$. A finite excursion expansion separates the constant direction from the compressed complement operator. The Hilbert--Schmidt bound leaves at most one exceptional eigenvalue. Positivity and boundedness of the kernel then force compensation through the two components of the degree fluctuation. A single truncated-square inequality proves that this compensation suffices at every cutoff; positive integration supplies all odd lengths $m\ge5$. The remaining scalar argument consists of fixed low-degree identities and completions of squares. The triangle is proved separately. Together with the case $p\ge2/3$, treated in a companion paper, this establishes the inequality at every edge density.
\end{abstract}

\paragraph{Context.}
The inequality $t(C_m,W)\ge p^m-p(1-p)^{m-1}$ was proved for $p\ge2/3$ in \cite{kim2026goodman}, where equality is attained by the balanced complete $k$-partite graphon at each density $p=1-1/k$. For $m=3$ the inequality is Goodman's triangle bound~\cite{goodman1959sets}, valid at every edge density. This article treats the remaining range. The excursion expansion of Section~\ref{sec:excursions} is the cycle case of \cite[Lemma~3.1]{kim2026goodman}, written in the notation used here; its proof is included so that the article is self-contained.

\section{The theorem and the elementary reductions}\label{sec:reductions}
For a finite simple graph $F$ and a graphon $W:[0,1]^2\to[0,1]$, symmetric and measurable, write the homomorphism density and the edge density~\cite{lovasz2006limits,lovasz2012large} as
\[
 t(F,W)=\int_{[0,1]^{V(F)}}\prod_{ij\in E(F)}W(x_i,x_j)\prod_{i\in V(F)}dx_i,
 \qquad p=\int_{[0,1]^2}W.
\]
\begin{theorem}\label{thm:main}
If $0\le p\le2/3$, then, for every odd integer $m\ge3$,
\begin{equation}\label{eq:target}
 t(C_m,W)\ge p^m-p(1-p)^{m-1}.
\end{equation}
\end{theorem}
The endpoint $p=2/3$ is also covered by \cite[Theorem~1.1]{kim2026goodman}; the argument below includes it as well.

\paragraph{Edge densities at most $1/2$.}
Put $q=1-p$. If $p\le1/2$, then $p^{m}-pq^{m-1}=p(p^{m-1}-q^{m-1})\le0$, whereas $t(C_m,W)\ge0$.

\paragraph{The triangle.}
For $m=3$, \eqref{eq:target} is Goodman's bound~\cite{goodman1959sets}; we include the standard short proof. Let $d_W(x)=\int_0^1 W(x,z)\,dz$. For $a,b\in[0,1]$, $ab\ge a+b-1$, so
\[
 \int W(x,z)W(y,z)\,dz\ge d_W(x)+d_W(y)-1.
\]
Multiplying by $W(x,y)$ and integrating yields
\[
 t(C_3,W)\ge2\int d_W(x)^2\,dx-p\ge2p^2-p=p^3-pq^2.
\]
The last inequality follows from $\int(d_W-p)^2\ge0$.

\paragraph{Reduction to step graphons.}
The following averaging argument is that of \cite[Lemma~2.1 and Appendix~A]{kim2026goodman}; it does not depend on the edge density. Let $W_N$ be the averages of $W$ over the rectangles of the level-$N$ dyadic partition. Then $W_N$ is a symmetric step graphon, $0\le W_N\le1$, and $\int W_N=p$ exactly. Moreover $W_N\to W$ in $L^1$: if $S$ is a dyadic rectangle step function with $\norm{W-S}_1<\varepsilon$, then for every sufficiently fine partition, averaging fixes $S$ and is an $L^1$ contraction, giving $\norm{W_N-W}_1<2\varepsilon$. Such step functions are dense in $L^1([0,1]^2)$.
Telescoping the product of the $m$ edge factors gives
\begin{equation}\label{eq:continuity}
 |t(C_m,W_N)-t(C_m,W)|\le m\norm{W_N-W}_1\longrightarrow0.
\end{equation}
It therefore suffices to prove the assertion for step graphons. From now on assume
\begin{equation}\label{eq:range}
 \frac12<p\le\frac23,\qquad \frac13\le q<\frac12,
\end{equation}
and work in the finite-dimensional space $E$ of functions constant on the cells of a partition for $W$ into finitely many cells of positive measure. For a kernel $K$ constant on the products of these cells, $T_K$ maps $L^2[0,1]$ into $E$ and vanishes on $E^\perp$; we regard $T_K$ as an operator on $E$, and all traces are taken on $E$. The space $E$ is closed under $f\mapsto|f|$ and $f\mapsto f_\pm$.

\section{The block decomposition and the spectral bound}
Put $U=1-W$ and let $T_U,T_W$ denote the corresponding integral operators. The function $\one$ has norm one. On $E=\R\one\oplus E_0$, where $E_0=E\cap\one^\perp$, define
\[
 g=T_U\one-q\one,\qquad A=\Pi_{E_0}T_U|_{E_0}.
\]
Then
\begin{equation}\label{eq:blocks}
 T_U=\begin{pmatrix}q&g^*\\g&A\end{pmatrix},\qquad
 T_W=\begin{pmatrix}p&-g^*\\-g&-A\end{pmatrix}.
\end{equation}
Here $g=p-d_W$ is the mean-zero degree fluctuation with its sign reversed, and $A$ is self-adjoint. This is the block form \cite[Eq.~(2.2)]{kim2026goodman}, there written for cells of equal measure.

If a cell $B_i$ has measure $b_i$, the matrix of a step-kernel operator $T_K$ in the orthonormal basis $\one_{B_i}/\sqrt{b_i}$ has entries $K_{ij}\sqrt{b_ib_j}$. Matrix multiplication around a cycle therefore gives the trace formula~\cite{lovasz2012large}, for $m\ge2$,
\begin{equation}\label{eq:tracecycle}
 \Tr(T_K^m)=\int_{[0,1]^m}K(x_1,x_2)K(x_2,x_3)\cdots K(x_m,x_1)\,dx_1\cdots dx_m .
\end{equation}
For $m\ge3$ the right side is $t(C_m,K)$; for $m=2$ it is $\int K^2$.

\begin{lemma}[Hilbert--Schmidt bound]\label{lem:HS}
One has
\begin{equation}\label{eq:HS}
 \Tr(A^2)+2\norm g^2\le pq.
\end{equation}
In particular every eigenvalue $\lambda$ of $A$ satisfies $|\lambda|\le\sqrt{pq}<p$.
\end{lemma}
\begin{proof}
Because $0\le U\le1$,
\[
 q^2+2\norm g^2+\Tr(A^2)=\Tr(T_U^2)=\int U^2\le\int U=q.
\]
The first equality is the block form \eqref{eq:blocks}, and the second is the trace formula \eqref{eq:tracecycle} with $m=2$. Subtract $q^2$. The eigenvalue assertion follows from $\lambda^2\le\Tr(A^2)$ and $q<p$.
\end{proof}

\section{Excursions: an exact finite counting identity}\label{sec:excursions}
This section supplies both the interpretation and the inequality used later. Insert
$I=\Pi_0+\Pi_H$, the projections onto $\R\one$ and $E_0$, between consecutive copies of $T_W$ in the trace. Thus
\begin{equation}\label{eq:words}
 \Tr(T_W^m)=\sum_{\omega\in\{0,H\}^m}
 \Tr\bigl(\Pi_{\omega_1}T_W\Pi_{\omega_2}T_W\cdots\Pi_{\omega_m}T_W\bigr).
\end{equation}
Indices are cyclic. An \emph{excursion} in a mixed word is a maximal nonempty cyclic block of $H$'s, with a $0$ before and after it. The $0$-only word contributes $p^m$; the $H$-only word contributes $\Tr((-A)^m)$.

An excursion with $j\ge0$ internal transitions in $E_0$, followed by $k\ge0$ constant-to-constant transitions, has length $j+k+2$ and weight
\[
 p^k\ip{g}{(-A)^j g}.
\]
The two minus signs on departure and return cancel. For $\ell\ge2$, define the sum of these weights at total length $\ell$ by
\begin{equation}\label{eq:e}
 e_\ell=\sum_{j=0}^{\ell-2}p^{\ell-2-j}\ip{g}{(-A)^j g}.
\end{equation}

\begin{lemma}[Cyclic excursion expansion]\label{lem:exc}
For every integer $m\ge3$,
\begin{equation}\label{eq:exactexc}
 \Tr(T_W^m)=p^m+\Tr((-A)^m)
 +\sum_{r=1}^{\lfloor m/2\rfloor}\frac mr
 \sum_{\substack{\ell_1+\cdots+\ell_r=m\\\ell_i\ge2}}
 e_{\ell_1}\cdots e_{\ell_r}.
\end{equation}
Every $e_\ell$ is nonnegative. Consequently, for odd $m\ge3$,
\begin{equation}\label{eq:oneexc}
 t(C_m,W)\ge p^m-\Tr(A^m)+m\ip{g}{k_m(A)g},
 \qquad k_m(z)=\frac{p^{m-1}-z^{m-1}}{p+z}.
\end{equation}
The function $k_m(A)$ can equivalently be defined by the polynomial
$\sum_{j=0}^{m-2}p^{m-2-j}(-A)^j$.
\end{lemma}
\begin{proof}
Consider the mixed words in \eqref{eq:words} having exactly $r$ excursions, and mark one of their departures from $0$ to $H$. Every such word has $r$ possible marks. Conversely, choose the position of the marked departure among the $m$ indexed positions, and then choose an ordered list of $r$ excursions, each together with the constant transitions following its return. These data uniquely determine the marked word. Since $\Pi_0$ has rank one, its weight is the product of the $r$ scalar block weights. Summing these weights gives
\[
 r\,\bigl(\text{total weight of words with $r$ excursions}\bigr)
 =m\sum_{\substack{\ell_1+\cdots+\ell_r=m\\\ell_i\ge2}}
 e_{\ell_1}\cdots e_{\ell_r}.
\]
This proves \eqref{eq:exactexc}. Marking indexed departures also handles periodic words; no assumption about the number of distinct cyclic rotations is made.

Choose an orthonormal eigenbasis $(\phi_i)$ of $A$, with eigenvalues $\lambda_i$, and put $g_i=\ip g{\phi_i}$. Summing the finite geometric series in \eqref{eq:e} gives
\begin{equation}\label{eq:epos}
 e_\ell=\sum_i g_i^2\frac{p^{\ell-1}-(-\lambda_i)^{\ell-1}}{p+\lambda_i}\ge0.
\end{equation}
The denominator and numerator are positive, since $|\lambda_i|<p$ by Lemma~\ref{lem:HS}. Thus every aggregate with $r\ge2$ in \eqref{eq:exactexc} is nonnegative. For odd $m$, the $r=1$ aggregate is $m e_m=m\ip g{k_m(A)g}$, and $\Tr((-A)^m)=-\Tr(A^m)$. Discarding $r\ge2$ proves \eqref{eq:oneexc}.
\end{proof}

\begin{remark}
Individual projection words may have negative weight. Nonnegativity is proved only after the excursion lengths and constant transitions have been summed into the coefficients $e_\ell$. It is these aggregates, not individual words, that are discarded.
\end{remark}

\section{Isolating the exceptional eigenvalue}\label{sec:isolate}
There is at most one eigenvalue of $A$ greater than $q$, counted with multiplicity: two such eigenvalues would give
\[
 \Tr(A^2)>2q^2\ge pq,
\]
contrary to \eqref{eq:HS}; the second inequality uses $q\ge1/3$.

If every eigenvalue is at most $q$, then for odd $m$,
\[
 \lambda^m\le q^{m-2}\lambda^2
\]
for every eigenvalue $\lambda$, including the negative ones. Hence
\[
 \Tr(A^m)\le q^{m-2}\Tr(A^2)\le pq^{m-1}.
\]
The last term of \eqref{eq:oneexc} is $m\ip{g}{k_m(A)g}=me_m\ge0$ by \eqref{eq:epos}, so \eqref{eq:oneexc} gives $t(C_m,W)\ge p^m-pq^{m-1}$.

It remains to consider a simple eigenvalue $\Lambda>q$. Choose a real unit eigenfunction $\phi$ and write
\begin{equation}\label{eq:split}
 A\phi=\Lambda\phi,\qquad
 g=\gamma\phi+g_s,\qquad g_s\perp\one,\phi,\qquad
 s=\norm{g_s},\qquad M=\sqrt{pq-\Lambda^2}.
\end{equation}
Let $A_s$ be the restriction of $A$ to $\{\one,\phi\}^\perp\cap E$. The Hilbert--Schmidt bound gives
\begin{equation}\label{eq:safeHS}
 \Tr(A_s^2)+2(\gamma^2+s^2)\le M^2.
\end{equation}
Every eigenvalue of $A_s$ therefore lies in $[-M,M]$, and
\begin{equation}\label{eq:ordering}
 0\le M<q<\Lambda<\frac12<p.
\end{equation}
Indeed $M^2<pq-q^2=q(1-2q)\le q^2$, while $\Lambda^2\le pq<1/4$.

\paragraph{The three-state excursion interpretation.}
In the decomposition $\R\one\oplus\R\phi\oplus\{\one,\phi\}^\perp$,
\begin{equation}\label{eq:threeblocks}
 T_W=\begin{pmatrix}
 p&-\gamma&-g_s^*\\
 -\gamma&-\Lambda&0\\
 -g_s&0&-A_s
 \end{pmatrix}.
\end{equation}
There is no transition between the exceptional direction and the remaining spectral subspace. Consequently a nonzero excursion stays in one of these two subspaces. An exceptional excursion with $b\ge1$ nonconstant positions and $m-b\ge1$ constant positions has, for each rotation, weight
\[
 p^{m-b-1}\gamma^2(-\Lambda)^{b-1}.
\]
Summing over its $m$ rotations and $b=1,\ldots,m-1$ gives $m k_m(\Lambda)\gamma^2$. The corresponding remaining-subspace sum is $m\ip{g_s}{k_m(A_s)g_s}$.

For $0\le z<p$, $k_m(-z)\ge k_m(z)$, and, with $n=m-1$,
\[
 k_m'(z)=-\frac{p^n+npz^{n-1}+(n-1)z^n}{(p+z)^2}<0.
\]
Thus $k_m(\lambda)\ge k_m(M)$ for $|\lambda|\le M$. Since $A=\Lambda\oplus A_s$ on $\R\phi\oplus(\{\one,\phi\}^\perp\cap E)$, the polynomial $k_m(A)$ splits in the same way, and
\[
 \ip{g}{k_m(A)g}=k_m(\Lambda)\gamma^2+\ip{g_s}{k_m(A_s)g_s}
 \ge k_m(\Lambda)\gamma^2+k_m(M)s^2,
\]
where the inequality expands $g_s$ in an eigenbasis of $A_s$. Hence
\begin{equation}\label{eq:couplelower}
 m\ip{g}{k_m(A)g}\ge m k_m(\Lambda)\gamma^2+m k_m(M)s^2.
\end{equation}
The words that never visit the constant direction contribute exactly $-\Lambda^m-\Tr(A_s^m)$. Since $m$ is odd, \eqref{eq:safeHS} gives
\begin{align}\label{eq:noreturn}
 \Tr(A^m)
 &\le\Lambda^m+M^{m-2}\Tr(A_s^2)\notag\\
 &\le\Lambda^m+M^m-2M^{m-2}(\gamma^2+s^2).
\end{align}
Combining \eqref{eq:oneexc}, \eqref{eq:couplelower}, and \eqref{eq:noreturn},
\begin{equation}\label{eq:Ebound}
 t(C_m,W)-(p^m-pq^{m-1})\ge E_m+2M^{m-2}(\gamma^2+s^2),
\end{equation}
where
\begin{align}\label{eq:E}
 E_m={}&\frac{m\gamma^2}{p+\Lambda}(p^{m-1}-\Lambda^{m-1})
 +\frac{ms^2}{p+M}(p^{m-1}-M^{m-1})\notag\\
 &-\Lambda^m-M^m+pq^{m-1}.
\end{align}
We will prove $E_m\ge0$ without using the extra nonnegative term in \eqref{eq:Ebound}.

\section{Compensation forced by the eigenfunction}
The next lemma eliminates the shape of the eigenfunction. Its coefficient is chosen to match the scalar inequality used later.
\begin{lemma}\label{lem:comp}
One has $M>0$. The sign of $\phi$ may be chosen so that
\begin{equation}\label{eq:comp}
 \Lambda-q\le c\gamma_++\frac{s^2}{\Lambda-M},
 \qquad c=\sqrt{\frac{M^2(p-\Lambda)}{\Lambda pq}},
 \qquad \gamma_+=\max\{\gamma,0\}.
\end{equation}
\end{lemma}
\begin{proof}
Choose the sign so that $b=\ip{|\phi|}{\phi}\ge0$, and write, using notation local to this proof (recall that $|\phi|,\phi_\pm\in E$),
\[
 |\phi|=a\one+b\phi+k,\qquad k\perp\one,\phi.
\]
Then $a=\int|\phi|>0$ and $a^2+b^2+\norm k^2=1$. Since $\int\phi=0$, the positive and negative parts of $\phi$ both have integral $a/2$. Dropping the nonpositive cross term in $\ip\phi{T_U\phi}$ and using $U\le1$, as in the proof of \cite[Lemma~2.2]{kim2026goodman}, gives
\begin{equation}\label{eq:alower}
 \Lambda\le\left(\int\phi_+\right)^2+\left(\int\phi_-\right)^2=\frac{a^2}{2}.
\end{equation}
Also $U\ge0$ implies
\begin{equation}\label{eq:abs}
 \ip{|\phi|}{T_U|\phi|}\ge\ip\phi{T_U\phi}=\Lambda.
\end{equation}

\paragraph{A constraint on the spectral parameters.}
Use the bound $A\le\Lambda I$ on $E_0$ in \eqref{eq:abs}. Cauchy--Schwarz gives
\[
 \Lambda\le qa^2+2a\norm g\sqrt{1-a^2}+\Lambda(1-a^2),
\]
so
\[
 (\Lambda-q)a\le2\norm g\sqrt{1-a^2}.
\]
In particular $a<1$ and $\norm g>0$. By \eqref{eq:alower},
\[
 \norm g^2\ge\frac{(\Lambda-q)^2a^2}{4(1-a^2)}
 \ge\frac{\Lambda(\Lambda-q)^2}{2(1-2\Lambda)}.
\]
Combining this with $2\norm g^2\le M^2$ shows first that $M>0$, and then that
\begin{align}\label{eq:admiss}
 0&\le M^2(1-2\Lambda)-\Lambda(\Lambda-q)^2\notag\\
  &=(p-\Lambda)\bigl(q(1-\Lambda)-\Lambda^2\bigr).
\end{align}
Hence $q(1-\Lambda)-\Lambda^2\ge0$.

\paragraph{Separating the two coupling components.}
Expand \eqref{eq:abs} and now use $\ip k{Ak}\le M\norm k^2$ to obtain
\[
 (\Lambda-q)a^2+(\Lambda-M)\norm k^2
 \le2ab\gamma+2a\ip{g_s}{k}.
\]
Dividing by $a^2$ and completing a square in $\norm k/a$ yields
\begin{equation}\label{eq:compintermediate}
 \Lambda-q\le\frac{2b\gamma}{a}+\frac{s^2}{\Lambda-M}.
\end{equation}
If $\gamma\le0$, \eqref{eq:comp} follows immediately because $b\ge0$.

Suppose $\gamma>0$. Pointwise, $T_U\phi\le\int\phi_+=a/2$. Since $T_U\phi=\gamma\one+\Lambda\phi$ and $\norm{\phi_+}^2=(1+b)/2$, taking the inner product with $\phi_+$ gives
\[
 2a\gamma+2\Lambda(1+b)\le a^2\le1-b^2.
\]
Thus $0\le b\le1-2\Lambda$. Since $b^2/(1+b)$ is increasing for $b\ge0$,
\[
 \frac{b^2}{a^2}\le\frac{b^2}{2\Lambda(1+b)}
 \le\frac{(1-2\Lambda)^2}{4\Lambda(1-\Lambda)}.
\]
Consequently \eqref{eq:compintermediate} gives the sharper estimate
\begin{equation}\label{eq:sharpercomp}
 \Lambda-q\le\frac{1-2\Lambda}{\sqrt{\Lambda(1-\Lambda)}}\gamma_+
 +\frac{s^2}{\Lambda-M}.
\end{equation}

\paragraph{Matching the coefficient to the spectral parameters.}
Write $\varepsilon=\Lambda-q$. The identity
\begin{align}\label{eq:parametercert}
 &(1-\Lambda)M^2(p-\Lambda)-pq(1-2\Lambda)^2\notag\\
 &\quad=\varepsilon\Bigl[(1-2\Lambda)(q-p\Lambda)
 +\varepsilon\bigl(q(1-\Lambda)-\Lambda^2\bigr)\Bigr]
 \ge0
\end{align}
follows by expansion with $p+q=1$ and $M^2=pq-\Lambda^2$. Its sign uses \eqref{eq:admiss} and
\[
 q-p\Lambda\ge q-\frac p2=\frac{3q-1}{2}\ge0.
\]
After division by $\Lambda pq(1-\Lambda)>0$, it states
\[
 \frac{(1-2\Lambda)^2}{\Lambda(1-\Lambda)}\le c^2.
\]
Thus \eqref{eq:sharpercomp} implies \eqref{eq:comp}.
\end{proof}

\section{Removing the cycle length by positive integration}
For $t,z\ge0$, define
\begin{equation}\label{eq:D}
 D_t(z)=\frac{d}{dz}\bigl(z^2(z-t)_+^2\bigr)
       =2z(z-t)_+(2z-t).
\end{equation}
Define the cutoff expression
\begin{align}\label{eq:K}
 \mathcal K(t)={}&
 \frac{\gamma^2}{p+\Lambda}\bigl(D_t(p)-D_t(\Lambda)\bigr)
 +\frac{s^2}{p+M}\bigl(D_t(p)-D_t(M)\bigr)\notag\\
 &+pq(q-t)_+^2-\Lambda^2(\Lambda-t)_+^2-M^2(M-t)_+^2.
\end{align}
Let $m\ge5$ be odd, put $k=m-5\ge0$ and $C_m=(m-2)(m-3)(m-4)/2$, and let $0\le z\le p$. Both integrands below vanish for $t>z$, so the integrals over $[0,p]$ are integrals over $[0,z]$. The beta integral
\[
 \int_0^z t^{k}(z-t)^2\,dt=\frac{2z^{k+3}}{(k+1)(k+2)(k+3)}
\]
gives
\begin{equation}\label{eq:powerint}
 z^{m-2}=C_m\int_0^p t^{m-5}(z-t)_+^2\,dt.
\end{equation}
Since $D_t(z)=2z(2z^2-3zt+t^2)$ for $t\le z$,
\begin{align}\label{eq:powerintD}
 \int_0^p t^{k}D_t(z)\,dt
 &=2z^{k+4}\left(\frac2{k+1}-\frac3{k+2}+\frac1{k+3}\right)\notag\\
 &=\frac{2(k+5)z^{k+4}}{(k+1)(k+2)(k+3)}
 =\frac{mz^{m-1}}{C_m}.
\end{align}
Apply \eqref{eq:powerintD} with $z=p,\Lambda,M$ to the three powers $p^{m-1},\Lambda^{m-1},M^{m-1}$ in \eqref{eq:E}, and \eqref{eq:powerint} to the remaining terms written as
\[
 \Lambda^m=\Lambda^2\cdot\Lambda^{m-2},\qquad
 M^m=M^2\cdot M^{m-2},\qquad
 pq^{m-1}=pq\cdot q^{m-2}.
\]
All four points $p,\Lambda,M,q$ lie in $[0,p]$ by \eqref{eq:ordering}. Comparing with \eqref{eq:K} term by term yields
\begin{equation}\label{eq:Eint}
 E_m=C_m\int_0^p t^{m-5}\mathcal K(t)\,dt
 \qquad(m\ge5).
\end{equation}
It therefore remains to prove $\mathcal K(t)\ge0$ for every cutoff $t\in[0,p]$. This is a stronger assertion than positivity of the integrals alone, and it contains no cycle-length parameter.

\section{The length-independent scalar lemma}
All notation in the next lemma is local. The exceptional spectral level is normalized to one.
\begin{lemma}[Truncated-square inequality]\label{lem:scalar}
Suppose
\begin{equation}\label{eq:normalized}
 0<y<x<1<P,\qquad Px=1+y^2,\qquad P\le2x.
\end{equation}
Set $h=1-x$, and let $u,v\ge0$ satisfy $u+v\ge h$. Put
\[
 A(t)=\frac{D_t(P)-D_t(1)}{P+1},\qquad
 B(t)=\frac{D_t(P)-D_t(y)}{P+y}.
\]
Then, for $0\le t\le P$,
\begin{align}\label{eq:Phi}
 \Phi(t)={}&\frac{1+y^2}{y^2(P-1)}A(t)u^2+(1-y)B(t)v\notag\\
 &+(1+y^2)(x-t)_+^2-(1-t)_+^2-y^2(y-t)_+^2
 \ge0.
\end{align}
\end{lemma}
\begin{proof}
For fixed $t\ge0$, $D_t(z)$ is nonnegative and nondecreasing in $z\ge0$. Indeed it vanishes for $z\le t$, and for $z\ge t$,
\[
 \frac{d}{dz}D_t(z)=12(z-t)^2+12t(z-t)+2t^2\ge0.
\]
In particular $A(t),B(t)\ge0$.

\subsection*{Two coefficient bounds}
The bounds needed above the cutoff $y$ are
\begin{align}
 A(t)&\ge P-1 &&(0\le t\le1),\label{eq:Abound}\\
 (1-y)B(t)&\ge2(1-t) &&(y\le t\le1).\label{eq:Bbound}
\end{align}
For the first, on $[0,1]$,
\[
 A'(t)=\frac{2(P-1)}{P+1}\bigl(2t-3(P+1)\bigr)<0.
\]
Thus
\[
 A(t)\ge A(1)=\frac{2P(P-1)(2P-1)}{P+1}\ge P-1;
\]
the last difference equals $(P-1)^2(4P+1)/(P+1)$.

For \eqref{eq:Bbound}, set $r=1-t$, so $0\le r\le1-y$. Here $D_t(y)=0$ and
\[
 B(t)=2\frac{P}{P+y}(P-t)(2P-t).
\]
Each displayed factor is nonnegative and nondecreasing in $P\ge1+y^2$, and $P\ge1+y^2$ follows from $Px=1+y^2$ and $x<1$. It suffices to substitute $P=1+y^2$. After clearing the positive denominator and dividing by two, the needed inequality is certified by
\begin{align}\label{eq:Bcert}
 &(1-y)(1+y^2)(r+y^2)(r+1+2y^2)-r(1+y+y^2)\notag\\
 &\quad=(1-y)(1+y^2)(r-y)^2\notag\\
 &\qquad\quad+y^2r\bigl((1-y)^2+y^2(2-3y+2y^2)\bigr)\notag\\
 &\qquad\quad+2y^4(1+y^2)(1-y-r)\ge0.
\end{align}
All terms are nonnegative because $r\le1-y$ and
$2-3y+2y^2=2(y-3/4)^2+7/8>0$.

\subsection*{Cutoffs at or above $y$}
For $y\le t\le x$, write $r=1-t$, so $x-t=r-h$. Apply \eqref{eq:Abound}--\eqref{eq:Bbound}, and then use $v\ge h-u$:
\begin{align}\label{eq:middlesquares}
 \Phi(t)&\ge\frac{1+y^2}{y^2}u^2+2r(h-u)+(1+y^2)(r-h)^2-r^2\notag\\
 &=\left(y(r-h)-\frac uy\right)^2+(u-h)^2\ge0.
\end{align}
For $x\le t\le1$, $0\le r\le h$ and the reference term vanishes. Since $(1+y^2)/y^2\ge1$, the same bounds give
\begin{align}\label{eq:uppersquare}
 \Phi(t)&\ge u^2+2r(h-u)-r^2\notag\\
 &=(u-r)^2+2r(h-r)\ge0.
\end{align}
For $1\le t\le P$, only the nonnegative coupling terms involving $D_t(P)$ remain. Thus $\Phi(t)\ge0$ for every $t\ge y$.

\subsection*{A consequence of the density restriction}
To treat $t\le y$, we first record a bound on $h$. From $P\le2x$ and $Px=1+y^2$,
\[
 y^2\le2x^2-1=1-4h+2h^2.
\]
In particular $x>1/\sqrt2>1/2$, and
\[
 (2x-1)^2-y^2\ge2(1-x)^2\ge0.
\]
Hence $y\le2x-1$, or $h\le(1-y)/2$. Substituting this in the preceding inequality,
\[
 4h\le1-y^2+2h^2
 \le(1-y)(1+y)+\frac{(1-y)^2}{2}
 =\frac{(1-y)(3+y)}2.
\]
Therefore
\begin{equation}\label{eq:hbound}
 h\le\frac{(1-y)(3+y)}8.
\end{equation}
Set, for the rest of this proof,
\[
 H=(1+y^2)h-y^2(1-y).
\]
Since $h>0$, \eqref{eq:hbound} implies
\begin{align}\label{eq:Hbound}
 \frac{H}{(1-y)h}
 &=\frac{1+y^2}{1-y}-\frac{y^2}{h}\notag\\
 &\le\frac{3+4y-y^2}{3+y}\notag\\
 &=\frac32-\frac{(1-y)(3-2y)}{2(3+y)}\le\frac32.
\end{align}

\subsection*{Monotonicity for cutoffs below $y$}
On $[0,y]$, all positive parts in \eqref{eq:Phi} can be removed. The resulting polynomial is a convex quadratic because
\begin{equation}\label{eq:convex}
 \Phi''(t)=\frac{4(1+y^2)}{y^2(P+1)}u^2
 +\frac{4(1-y)(P-y)}{P+y}v\ge0.
\end{equation}
Its left derivative at $y$ is
\begin{align}\label{eq:derivative}
 \frac12\Phi'(y^-)=H
 &-\frac{1+y^2}{y^2}\left(3-\frac{2y}{P+1}\right)u^2\notag\\
 &-\frac{(P-y)(3P+y)}{P+y}(1-y)v.
\end{align}
We claim that this is nonpositive. If $H\le0$, the claim is immediate. Suppose $H>0$, so
\begin{equation}\label{eq:Hpositive}
 h>\frac{y^2(1-y)}{1+y^2}.
\end{equation}
First, since $P>1>y$,
\begin{equation}\label{eq:firstcost}
 \frac{1+y^2}{y^2}\left(3-\frac{2y}{P+1}\right)
 \ge\frac{2(1+y^2)}{y^2}\ge\frac{2(1-y)}h.
\end{equation}
For the other coefficient, \eqref{eq:Hpositive} gives
\[
 P=\frac{1+y^2}{1-h}\ge(1+y^2)(1+h)
 \ge1+2y^2-y^3=:P_0.
\]
The function
\[
 \frac{(P-y)(3P+y)}{P+y}=3P-5y+\frac{4y^2}{P+y}
\]
increases in $P\ge y$: its derivative is $3-4y^2/(P+y)^2\ge2$. Here $P_0>y$. At $P=P_0$, the following identity proves that this function is at least two:
\begin{align}\label{eq:derivcert}
 &(P_0-y)(3P_0+y)-2(P_0+y)\notag\\
 &\quad=(1-2y)^2+y^2(1-y)(6y^2-5y+3)
 +3y^4(1-y)^2\ge0,
\end{align}
since $6y^2-5y+3=6(y-5/12)^2+47/24>0$. We have shown
\begin{equation}\label{eq:secondcost}
 \frac{(P-y)(3P+y)}{P+y}\ge2\qquad\text{whenever }H>0.
\end{equation}
By \eqref{eq:firstcost}--\eqref{eq:secondcost} and $v\ge h-u$, the sum of the two subtracted terms in \eqref{eq:derivative} is at least
\begin{align}\label{eq:finalsquare}
 \frac{2(1-y)}h u^2+2(1-y)v
 &\ge\frac{2(1-y)}h u^2+2(1-y)(h-u)\notag\\
 &=2(1-y)h\left[\left(\frac uh-\frac12\right)^2+\frac34\right]\notag\\
 &\ge\frac32(1-y)h\ge H,
\end{align}
where the final inequality is \eqref{eq:Hbound}. Thus $\Phi'(y^-)\le0$.

Convexity on the lower interval now gives $\Phi'(t)\le\Phi'(y^-)\le0$ for $0\le t<y$. Consequently $\Phi(t)\ge\Phi(y)\ge0$ for $t\le y$, completing the proof. Only the derivative of the polynomial on the lower interval is used; differentiability across the cutoff junction is not assumed.
\end{proof}

\section{Application of the scalar lemma and completion}
We verify all hypotheses of Lemma~\ref{lem:scalar} for the graphon parameters. Put
\begin{equation}\label{eq:normalizeapplication}
 P=\frac p\Lambda,\qquad x=\frac q\Lambda,\qquad
 y=\frac M\Lambda,\qquad h=1-x,
\end{equation}
and, with $c$ and the sign of $\phi$ as in Lemma~\ref{lem:comp}, put
\begin{equation}\label{eq:normalizecoupling}
 u=\frac{c\gamma_+}{\Lambda},\qquad
 v=\frac{s^2}{\Lambda(\Lambda-M)}.
\end{equation}
By \eqref{eq:ordering} and $M>0$,
\[
 0<y<x<1<P.
\]
Moreover $M^2+\Lambda^2=pq$ gives $Px=1+y^2$, while $q\ge1/3$ gives $p\le2q$, hence $P\le2x$. Inequality \eqref{eq:comp} is exactly $u+v\ge h$. Replacing $\phi$ by $-\phi$ changes $\gamma$ into $-\gamma$ and leaves $g_s$, $s$, and $E_m$ unchanged, since $E_m$ depends on $\gamma$ only through $\gamma^2$; so the choice of sign in Lemma~\ref{lem:comp} does not affect \eqref{eq:Ebound}.

The definition of $c$ yields
\[
 c^2=\frac{y^2(P-1)}{1+y^2},\qquad
 \frac{\gamma^2}{\Lambda^2}\ge
 \frac{1+y^2}{y^2(P-1)}u^2,\qquad
 \frac{s^2}{\Lambda^2}=(1-y)v.
\]
Finally $D_{\Lambda t}(\Lambda z)=\Lambda^3D_t(z)$, so \eqref{eq:K} and nonnegativity of its coupling coefficients give
\begin{align*}
 \frac{\mathcal K(\Lambda t)}{\Lambda^4}
 ={}&\frac{\gamma^2}{\Lambda^2}A(t)+\frac{s^2}{\Lambda^2}B(t)\notag\\
 &+(1+y^2)(x-t)_+^2-(1-t)_+^2-y^2(y-t)_+^2
 \ge\Phi(t)\ge0
\end{align*}
for every $t\in[0,P]$. Thus $\mathcal K\ge0$ throughout $[0,p]$.

For each odd $m\ge5$, \eqref{eq:Eint} gives $E_m\ge0$, and \eqref{eq:Ebound} proves
\[
 t(C_m,W)\ge p^m-pq^{m-1}.
\]
This settles the exceptional-eigenvalue case; the case without an exceptional eigenvalue was already proved in Section~\ref{sec:isolate}. The step approximation \eqref{eq:continuity} returns the result to arbitrary graphons. Together with the triangle argument and the case $p\le1/2$ in Section~\ref{sec:reductions}, this proves Theorem~\ref{thm:main}.\hfill$\square$

\section*{Structure and interpretation}
The proof has four substantive components. The excursion identity separates the constant contribution, the nonconstant trace, and the coupling contributions, with positivity justified only after the appropriate sums. The Hilbert--Schmidt bound leaves at most one compression eigenvalue above $q$. The absolute-value and positive-part comparisons of its eigenfunction force sufficient coupling either directly to that eigenfunction or to the remaining spectral subspace. The truncated-square lemma then compares these two contributions with the spectral deficit at every cutoff. All higher odd powers follow from the same positive integral representation.

Beyond completions of squares, the proof uses two polynomial certificates, \eqref{eq:Bcert} and \eqref{eq:derivcert}. The first lower-bounds the remaining-subspace excursion coefficient, and the second controls the derivative below the smallest cutoff. The central interval is settled by the two squares in \eqref{eq:middlesquares}. After the excursion reduction, no estimate depends on $m$.

\paragraph{Verification note.}
The accompanying file \nolinkurl{verify_complete_odd_cycle_proof.py} checks the displayed algebraic certificates, the normalization, representative exact excursion identities, and the power-integral formula using rational symbolic arithmetic. These checks supplement the mathematical proof; they are not a formalization in a proof assistant.

\bibliographystyle{abbrv}
\bibliography{references}
\end{document}
